\begin{abstract}
Tensegrity-inspired structures of rigid and soft members could cushion
impacts, but the shock they transmit is hard to predict by simulation.
We report a Bayesian optimization campaign, driven by drop tests, over
three-strut prisms printed in one piece from polylactic acid (PLA) and
thermoplastic polyurethane (TPU). The search space begins with five
geometry variables and, from the third batch, adds infill densities and
filament settings. Each candidate is scaled to a common mass,
3D printed, and drop-tested repeatedly. The optimizer seeks low
transmissibility (measured as a ratio of filtered top-to-base peak
accelerations) and low mass-weighted rebound (specimen mass times the
fraction of impact velocity returned at the top). Most 3D-printed
specimens amplify the peak shock. Over a quasi-random initial batch of
experiments and three model-recommended batches (35 specimens), the
dominated hypervolume grew 65 percent and the model chose every design
on the final Pareto front. The final batch's best-predicted design
lowered transmissibility to 0.803, ten percent below the best
initial-batch specimen. We printed one batch three times to measure
fabrication repeatability. Print-to-print scatter in transmissibility
was near one percent, and a seating check flagged both large outliers,
so we treat single-print extremes as provisional until replicated.
\end{abstract}
